% Part:sets-functions-relations
% Chapter: size-of-sets
% Section: pairing

\documentclass[../../../include/open-logic-section]{subfiles}

\begin{document}

\olfileid[tr]{sfr}{siz}{pai}
\olsection{İkileme Fonksiyonları ve Kodlar}

\begin{explain}
Cantor'un zikzak yöntemi $\Nat^n$'nin sayılabilirliğini görsel olarak açık
hâle getirir. Fakat $\Nat^2$'yi betimleyen dizimize odaklanalım. Dizideki
zikzak çizgisini izleyip konumları sayarak $\tuple{1,2}$'nin $7$ sayısıyla
ilişkili olduğunu denetleyebiliriz. Yine de bunu daha doğrudan hesaplamak iyi
olurdu. Başka bir deyişle, zikzak numaralandırmasının
\emph{tersi}ni, yani
\[
g(\tuple{0,0}) = 0, \;
g(\tuple{0,1}) = 1, \;
g(\tuple{1,0}) = 2, \; \dots,
g(\tuple{1,2}) = 7, \; \dots
\]
olacak biçimde bir $g\colon \Nat^2 \to \Nat$ fonksiyonunu elimizin altında
bulundurmak iyi olurdu. Bu, $\tuple{n,m}$'nin numaralandırmamızda tam olarak
nerede yer alacağını hesaplamamızı sağlardı.

Aslında iki gözlem yaparak $g$'yi doğrudan tanımlayabiliriz. Birincisi,
$n$'inci satır ile $m$'inci sütun $v$ değerini içeriyorsa $(n+1)$'inci satır
ile $(m-1)$'inci sütun $v+1$ değerini içerir. İkincisi,
numaralandırmamızın ilk satırı $0$, $1$, $3$, $6$ vb. ile başlayan üçgensel
sayılardan oluşur. $k$'inci üçgensel sayı, $k$'dan büyük olmayan doğal
sayıların toplamıdır ve $k(k+1)/2$ olarak hesaplanabilir. Bu iki gözlemi birleştirip şu
fonksiyonu düşünün:
\[
  g(n,m) = \frac{(n+m+1)(n+m)}{2} + n.
\]
Göze daha kolay geldiği için çoğu zaman $g(\tuple{n,m})$ yerine yalnızca
$g(n,m)$ yazarız. Bu formül önce $(n+m)$'inci üçgensel sayıyı belirlemenizi,
sonra buna $n$'yi eklemenizi söyler. Diziyi tam istediğimiz biçimde doldurur.
Özellikle $\tuple{1,2}$ ikilisi $\frac{4 \times 3}{2}+1=7$'ye gönderilir.

Bu $g$ fonksiyonu, bir ikililer kümesinin numaralandırmasının
\emph{tersi}dir. Böyle fonksiyonlara \emph{ikileme fonksiyonları} denir.
\end{explain}

\begin{defn}[İkileme fonksiyonu]
  Bir $f\colon A \times B \to \Nat$ fonksiyonu bire bir ise aritmetik bir
  \emph{ikileme fonksiyonu}dur. Ayrıca $f$'nin $A \times B$'yi
  \emph{kodladığını}, $f(x,y)$'nin de $\tuple{x,y}$'nin \emph{kodu}
  olduğunu söyleriz.
\end{defn}

\begin{explain}
İkileme fonksiyonlarını örneğin doğal sayı ikililerini kodlamak için
kullanabiliriz; başka bir deyişle her eleman \emph{ikilisini} \emph{tek} bir
sayıyla temsil edebiliriz. İkileme fonksiyonunun tersini kullanarak sayının
\emph{kodunu çözebilir}, yani hangi ikiliyi temsil ettiğini bulabiliriz.
\end{explain}

\begin{prob}
Negatif olmayan bütün rasyonel sayıların bir numaralandırmasını verin.
\end{prob}

\begin{prob}
$\Rat$'ın !!{enumerable} olduğunu gösterin. Her rasyonel sayının $z\in\Int$
ve $m\in\Nat^+$ olmak üzere $z/m$ kesri biçiminde yazılabileceğini hatırlayın.
\end{prob}

\begin{prob}
$\Bin^*$'ın bir numaralandırmasını tanımlayın.
\end{prob}

\begin{prob}
Giriş mantık dersinizden, her olası doğruluk tablosunun bir doğruluk
fonksiyonunu ifade ettiğini hatırlayın. Başka bir deyişle doğruluk
fonksiyonları, bazı $k$'lar için $\Bin^k \to \Bin$ biçimindeki bütün
fonksiyonlardır. Bütün doğruluk fonksiyonları kümesinin sayılabilir olduğunu
ispatlayın.
\end{prob}

\begin{prob}
Herhangi bir sonsuz !!{enumerable} kümenin bütün sonlu alt kümelerinden
oluşan kümenin !!{enumerable} olduğunu gösterin.
\end{prob}

\begin{prob}
$\Nat$'ın bir alt kümesinin \emph{tümleyeni sonlu} olması, $\Nat$ içindeki
tümleyeninin sonlu olması demektir; yani $A\subseteq\Nat$ için bu koşul
$\Nat\setminus A$'nın sonlu olmasıdır. Bir $I$ kümesinin !!{element}ları tam
olarak $\Nat$'ın sonlu ve tümleyeni sonlu alt kümeleri olsun. $I$'nın !!{enumerable}
olduğunu gösterin.
\end{prob}

\begin{prob}
!!{enumerable} kümelerin !!{enumerable} birleşiminin !!{enumerable} olduğunu
gösterin. Yani $A_1$, $A_2$, \dots{} kümelerinin her biri !!{enumerable}
olduğunda bunların tümünün birleşimi $\bigcup_{i=1}^\infty A_i$'nın da
!!{enumerable} olduğunu gösterin. [Not: Bu zordur!]
\end{prob}

\begin{editorial}
Bu sonuç, numaralandırmalar yalnızca varlık önermeleri olarak verildiğinde,
sayılabilir bir seçim ilkesine dayanır. Her $A_i$ için belirli bir
numaralandırma veri olarak sağlanmışsa ek bir seçim adımı gerekmez.
\end{editorial}

\begin{prob}
$f\colon A\times B\to\Nat$ herhangi bir ikileme fonksiyonu olsun.
$A\times B=\emptyset$ ise tanım gereği !!{enumerable}dir. $A\times B\neq
\emptyset$ ise $f$'nin görüntü kümesinin bir numaralandırmasını $f$'nin
tersiyle bileştirmenin $A\times B$'nin bir numaralandırmasını verdiğini gösterin.
\end{prob}

\begin{prob}
$\Nat^3$'ü kodlayan bir fonksiyon belirtin.
\end{prob}

\end{document}
